\chapter{From Mechanics to Fields}

\begin{remarkbox}
This chapter explains how the finite-dimensional language of classical
mechanics becomes the infinite-dimensional language of fields. The main idea is
simple but profound: a field is not one coordinate depending on time, but a
continuous family of coordinates, one for each point of space.
\end{remarkbox}

\section{Generalized Coordinates in Classical Mechanics}

Classical mechanics begins with the motion of bodies. In its most elementary
form, one describes a particle by its position as a function of time,
\[
    \vect{x}(t) = (x^1(t),x^2(t),x^3(t)).
\]
This description is already a function: the position is assigned to every time
at which the motion is considered. However, the number of independent
configuration variables at each time is finite. A single particle in three
spatial dimensions has three coordinates. A system of \(N\) particles has
\(3N\) coordinates.

The language of generalized coordinates abstracts away from the particular
choice of Cartesian coordinates. Instead of writing all particle positions
explicitly, one introduces coordinates
\[
    q^1(t),q^2(t),\ldots,q^n(t),
\]
where \(n\) is the number of degrees of freedom of the mechanical system. These
coordinates may be distances, angles, normal mode amplitudes, or any other
variables that describe the configuration of the system.

The important point is that the configuration at a fixed time is represented by
a point
\[
    q(t) = (q^1(t),\ldots,q^n(t))
\]
in an \(n\)-dimensional configuration space. The motion of the system is then a
curve in this configuration space.

\begin{definitionbox}
The \emph{configuration space} of a mechanical system is the space of all
possible instantaneous configurations of the system. If the system has
\(n\) generalized coordinates, a configuration is represented locally by a point
\((q^1,\ldots,q^n)\).
\end{definitionbox}

In the Lagrangian formulation, the dynamics is encoded in a Lagrangian
\[
    L(q,\dot q,t),
\]
where \(\dot q^i=dq^i/dt\). The action is the functional
\[
    S[q] = \int_{t_1}^{t_2} L(q(t),\dot q(t),t)\,dt.
\]
The physical trajectories are selected by the variational principle
\[
    \delta S = 0,
\]
for variations of the path that keep the endpoints fixed. This gives the
Euler--Lagrange equations
\[
    \frac{d}{dt}\left(\frac{\partial L}{\partial \dot q^i}\right)
    - \frac{\partial L}{\partial q^i}=0,
    \qquad i=1,\ldots,n.
\]

This structure will reappear almost unchanged in field theory. The difference
is that the finite list of coordinates \(q^i(t)\) will be replaced by fields
such as \(\phi(t,\vect{x})\). The ordinary derivative with respect to time will
be supplemented by spatial derivatives, and the ordinary Lagrangian \(L\) will
be replaced by a Lagrangian density \(\mathcal{L}\).

\section{Systems with Finitely Many Degrees of Freedom}

A system has finitely many degrees of freedom when its instantaneous state can
be described by finitely many independent numbers. A double pendulum, a finite
chain of masses and springs, a rigid body, or a finite electrical circuit are
examples of such systems.

For a system of \(N\) coupled oscillators arranged in a line, one may introduce
coordinates
\[
    q_1(t),q_2(t),\ldots,q_N(t),
\]
where \(q_i(t)\) is the displacement of the \(i\)-th oscillator from equilibrium.
A typical Lagrangian has the form
\[
    L = \sum_{i=1}^{N} \frac{m}{2}\dot q_i^2
        - \sum_{i=1}^{N} \frac{k}{2}(q_{i+1}-q_i)^2,
\]
with suitable boundary conventions. The first term describes kinetic energy;
the second describes elastic energy stored in relative displacements.

This example already contains the seed of field theory. The label \(i\) tells us
where the oscillator is located in the chain. If the oscillators are separated
by a distance \(a\), the position of the \(i\)-th oscillator is
\[
    x_i = i a.
\]
The coordinate \(q_i(t)\) may therefore be viewed as the value of a discrete
field at the lattice point \(x_i\):
\[
    q_i(t) \approx \phi(t,x_i).
\]

\begin{examplebox}
A finite chain of masses is not yet a field theory, because it has only
finitely many degrees of freedom. Nevertheless, it is the most useful bridge to
field theory: it shows how a spatially distributed system may be described by
one dynamical variable at each position.
\end{examplebox}

The equations of motion of a finite chain are ordinary differential equations.
For example, away from the endpoints one obtains equations of the schematic
form
\[
    m\ddot q_i = k(q_{i+1}-2q_i+q_{i-1}).
\]
The expression on the right-hand side is a discrete second derivative. This is
why the continuum limit of such a system naturally produces partial
differential equations.

\section{The Continuum Limit}

The continuum limit is the passage from a very large number of discrete degrees
of freedom to a field with continuously many degrees of freedom. The basic idea
is to let the spacing \(a\) between neighbouring points become small while the
number of points becomes large.

For the chain of oscillators, we introduce a field \(\phi(t,x)\) such that
\[
    q_i(t) \approx \phi(t,x_i),
    \qquad x_i = ia.
\]
Then the finite difference
\[
    \frac{q_{i+1}-q_i}{a}
\]
approximates the spatial derivative
\[
    \partial_x \phi(t,x).
\]
Similarly,
\[
    \frac{q_{i+1}-2q_i+q_{i-1}}{a^2}
\]
approximates the second spatial derivative
\[
    \partial_x^2 \phi(t,x).
\]
Thus an equation of motion involving nearest-neighbour couplings becomes, in an
appropriate limit, a differential equation for \(\phi(t,x)\).

The passage from sums to integrals is equally important. A discrete sum over
oscillators becomes an integral over space:
\[
    \sum_i a\, f(x_i) \longrightarrow \int f(x)\,dx.
\]
In mechanics, the action has the form
\[
    S[q] = \int L\,dt.
\]
In field theory, the action has the form
\[
    S[\phi] = \int \mathcal{L}(\phi,\partial_t\phi,\partial_x\phi,t,x)\,dt\,dx.
\]
In higher dimensions this becomes
\[
    S[\phi] = \int \mathcal{L}(\phi,\partial_\mu\phi,x)\,d^{d+1}x,
\]
where \(d\) is the number of spatial dimensions and \(d+1\) is the dimension of
spacetime.

\begin{remarkbox}
The continuum limit is not merely the replacement of sums by integrals. It also
changes the type of mathematics. Finite-dimensional configuration spaces are
replaced by spaces of functions, and ordinary differential equations are
replaced by partial differential equations.
\end{remarkbox}

A simple example is the wave equation. Starting from a chain of coupled
oscillators and taking the continuum limit, one obtains an equation of the form
\[
    \partial_t^2 \phi - c^2\partial_x^2\phi = 0,
\]
where \(c\) is the wave speed. This equation does not describe the motion of one
particle. It describes the evolution of a whole spatial profile.

\section{Fields as Infinite-Dimensional Dynamical Systems}

A field assigns a value to every point of space, or more generally to every
point of spacetime. For a real scalar field in space, a configuration at time
\(t\) is a function
\[
    \phi_t : \vect{x} \mapsto \phi(t,\vect{x}).
\]
Thus the instantaneous configuration is not a finite tuple of numbers, but an
entire function.

This is the precise sense in which field theory is infinite-dimensional. In a
mechanical system with coordinates \(q^i(t)\), the index \(i\) runs over a finite
set. In a field theory, the role of the index is partly played by the spatial
point \(\vect{x}\), which ranges over a continuum.

One may write the analogy informally as
\[
    q^i(t) \quad \longrightarrow \quad \phi(t,\vect{x}).
\]
The discrete label \(i\) is replaced by a continuous label \(\vect{x}\). The
velocity \(\dot q^i(t)\) becomes the time derivative \(\partial_t\phi(t,\vect{x})\).
The interaction between neighbouring degrees of freedom produces spatial
derivatives such as \(\nabla\phi\) or \(\Delta\phi\).

\begin{definitionbox}
A \emph{classical field} is a classical dynamical variable whose value depends
on spacetime position. A field configuration at a fixed time is a function on
space, and the time evolution of the field is the evolution of that function.
\end{definitionbox}

The configuration space of a field theory is therefore a space of functions. For
example, the configuration space of a real scalar field may be thought of as a
suitable space of functions
\[
    \phi : \mathbb{R}^d \to \mathbb{R}.
\]
The word ``suitable'' hides analytic details: one may need fields to be smooth,
square-integrable, rapidly decreasing, periodic, or compatible with boundary
conditions. These details matter in advanced treatments. At this stage, the
main point is conceptual: a field theory is a dynamical system whose
configuration space is infinite-dimensional.

\section{Examples: String, Membrane, Fluid Density, Electromagnetic Field}

Field theory is not limited to fundamental particle physics. It appears whenever
one studies quantities distributed over space and changing in time.

\subsection*{Vibrating string}

A vibrating string may be described by a displacement field
\[
    y(t,x),
\]
where \(x\) labels the position along the string and \(y\) is the transverse
displacement. The field equation is, in the simplest approximation,
\[
    \partial_t^2 y - c^2\partial_x^2 y = 0.
\]
The physical meaning is direct: the acceleration of a small piece of the string
is determined by the curvature of the string profile.

\subsection*{Membrane}

A membrane has two spatial labels. Its displacement may be written as
\[
    z(t,x,y),
\]
and its dynamics may involve equations such as
\[
    \partial_t^2 z - c^2(\partial_x^2 z + \partial_y^2 z)=0.
\]
The field now depends on time and two spatial coordinates.

\subsection*{Fluid density}

A fluid may be described by fields such as mass density \(\rho(t,\vect{x})\),
velocity \(\vect{v}(t,\vect{x})\), and pressure \(p(t,\vect{x})\). These fields
are not optional decorations; they are the natural variables because a fluid has
properties at each point in space.

The continuity equation
\[
    \partial_t \rho + \nabla\cdot(\rho\vect{v}) = 0
\]
expresses local conservation of mass. It says that density changes in a region
when mass flows into or out of that region.

\subsection*{Electromagnetic field}

Electromagnetism is the central example of a fundamental classical field theory.
The electric and magnetic fields
\[
    \vect{E}(t,\vect{x}), \qquad \vect{B}(t,\vect{x})
\]
assign vectors to every point of space and time. Maxwell's equations determine
how these fields are produced by charges and currents and how they propagate.

In relativistic notation, the electromagnetic field is described by the
antisymmetric field strength tensor \(F_{\mu\nu}\), or equivalently by a gauge
potential \(A_\mu\). This course will return to electromagnetism as the first
major example of a relativistic gauge field theory.

\begin{examplebox}
The same mathematical pattern appears in very different physical systems:
strings, membranes, fluids, elastic media, electromagnetic fields, and gauge
fields. The details differ, but the structural idea is the same: the dynamical
variables are functions on space or spacetime.
\end{examplebox}

\section{Initial Data and Boundary Data}

A field equation usually has many solutions. To select a particular solution,
one must specify appropriate data. For evolution equations, the most common data
are initial data: the field and enough time derivatives at an initial time.

For the wave equation
\[
    \partial_t^2\phi - c^2\Delta\phi = 0,
\]
one typically specifies
\[
    \phi(0,\vect{x}) = f(\vect{x}),
    \qquad
    \partial_t\phi(0,\vect{x}) = g(\vect{x}).
\]
The first function gives the initial profile. The second gives the initial
velocity profile. Together they determine the later evolution, under suitable
conditions.

Boundary data are also essential when the spatial domain has a boundary. For a
field on an interval \([0,L]\), one may impose fixed endpoint conditions
\[
    \phi(t,0)=0,
    \qquad
    \phi(t,L)=0,
\]
or periodic boundary conditions
\[
    \phi(t,0)=\phi(t,L).
\]
Different boundary conditions describe different physical systems and lead to
different spectra of allowed modes.

\begin{remarkbox}
Initial conditions select a history from many possible histories. Boundary
conditions define the physical problem itself: they say what kind of system the
field lives on and how it interacts with the edge of its domain.
\end{remarkbox}

Boundary terms will become especially important in the variational principle.
When deriving field equations from an action, integrations by parts produce
terms on the boundary. A variational principle is well-defined only when those
terms vanish or are cancelled by appropriate boundary contributions.

\section{Locality and Causality}

Most classical field theories used in physics are local. Locality means that the
equations of motion at a point involve the field and finitely many of its
derivatives at that same point, not the values of the field everywhere else.
For example, the wave equation
\[
    \partial_t^2\phi - c^2\Delta\phi = 0
\]
is local because it relates derivatives of \(\phi\) at the same spacetime point.

A nonlocal equation would involve expressions such as
\[
    \int K(\vect{x},\vect{y})\phi(t,\vect{y})\,d^d y,
\]
where the value of the equation at \(\vect{x}\) depends directly on field values
at many other points \(\vect{y}\). Nonlocal models exist and can be useful, but
locality is the organizing principle of standard classical field theory.

\begin{definitionbox}
A field equation is \emph{local} if the equation at a spacetime point depends
only on the values of the fields and a finite number of their derivatives at
that same point.
\end{definitionbox}

In relativistic theories, locality is closely connected with causality. Signals
and disturbances should not propagate faster than light. The causal structure of
Minkowski spacetime divides directions into timelike, lightlike, and spacelike.
A relativistic field equation should respect this structure.

The wave equation illustrates the idea. If a disturbance is created at one
point, it propagates with finite speed. In relativistic theories this speed is
bounded by the speed of light. This requirement is one reason why relativistic
field theories are naturally formulated on spacetime rather than on space alone.

\section{What Makes a Theory a Field Theory?}

A theory is a field theory when its fundamental dynamical variables are fields
and its equations determine how those fields vary over space and time. The
following features are typical.

\begin{enumerate}
    \item The configuration at a fixed time is a function, or a collection of
    functions, on space.
    \item The action is usually an integral of a Lagrangian density over
    spacetime.
    \item The equations of motion are usually partial differential equations.
    \item The theory has local degrees of freedom, meaning that different
    regions of space can carry different field values.
    \item Boundary conditions and initial data are part of the mathematical and
    physical specification of the problem.
\end{enumerate}

The transition from mechanics to field theory may therefore be summarized as
\[
    \text{finite coordinates } q^i(t)
    \quad \longrightarrow \quad
    \text{fields } \phi^A(t,\vect{x}),
\]
where \(A\) labels the type or component of the field. For example, a scalar
field has one component, a vector field has several components, and a gauge
field carries spacetime and possibly internal indices.

\begin{remarkbox}
The slogan ``a field is infinitely many coupled oscillators'' is useful but not
complete. It captures the continuum of degrees of freedom, but a full field
theory also includes locality, boundary conditions, symmetries, conservation
laws, and often relativistic spacetime structure.
\end{remarkbox}

This chapter has introduced the conceptual bridge from mechanics to fields. The
next chapter develops the mathematical language needed to work with fields:
indices, tensors, derivatives, differential forms, variations, distributions,
and spacetime notation.
